First suppose $L/K$ is separable. Embed $L$ in a finite Galois extension of $K$, and let $C$ be the integral closure of $A$ there. By \exref{5.13} and \exref{5.14}, there are finitely many primes of $C$ above any given prime of $A$. Since $C$ is integral over $B$, every prime of $B$ lifts to one of $C$. Thus the fibers for $B$ are finite.

In characteristic $p>0$, suppose $L/K$ is purely inseparable. Every $b\in B$ has $b^{p^n}\in K$ for some $n\geq0$; this power is integral over $A$, so it belongs to $A$. A prime $\mathfrak{q}$ above $\mathfrak{p}$ therefore satisfies
\[\mathfrak{q}=\{b\in B:b^{p^n}\in\mathfrak{p}\text{ for some }n\geq0\}.\]
Indeed, either membership implies the other by contraction and primality. Hence the prime is unique, and lying-over gives existence.

For general $L/K$, let $E$ be the maximal separable subextension and $D$ the integral closure of $A$ in $E$. Then $E/K$ is finite separable and $L/E$ is purely inseparable. Each $b\in B$ has a $p$-power in $E$, and that power lies in $D$ because it is integral over $A$. The same argument gives exactly one prime of $B$ above each prime of $D$. The first case gives finitely many primes of $D$ above each prime of $A$, proving the result. In characteristic zero, the separable case already applies.
